\section{光谱合成 Spectral Synthesis}\label{sec:synthesis}

光谱合成把模型大气和谱线表变成出射光谱。它计算贯穿大气的连续与谱线不透明度，然后求解辐射转移以确定离开恒星表面的辐射。\paynezero\ 用一种新的波长并行计算实现了其所支持的 Kurucz 合成物理 \citep{Kurucz2005}。

\begin{EN}
Spectral synthesis turns a model atmosphere and a line list into an emergent
spectrum. It calculates continuous and line opacity through the atmosphere,
then solves radiative transfer to determine the radiation that leaves the
stellar surface. \paynezero\ implements the supported
Kurucz synthesis physics \citep{Kurucz2005} in a new
wavelength-parallel calculation.
\end{EN}

深度状态一旦固定，不同波长之间的相互依赖就微乎其微。因此，宽波段光谱可以把一个长长的波长循环替换成同一短深度计算的许多副本。正是这种并行结构，使下文中的宽波段 Kurucz 合成基准从十余分钟降到数秒——而无需任何模拟器。我们先讲合成，因为它最直接地暴露了本文的核心计算思想。在展开计算之前，表~\ref{tab:atmosphere_symbols} 汇总了全文使用的符号。

\begin{EN}
Once the depth state is fixed, different wavelengths have little dependence on
one another. A broad spectrum can therefore replace one long wavelength loop
with many copies of the same short depth calculation. This parallel structure
is why the broad Kurucz synthesis benchmark below moves from more than ten
minutes to seconds without an emulator. We begin with synthesis because it exposes the central
computational idea directly.
Before developing that calculation, Table~\ref{tab:atmosphere_symbols}
collects the notation used throughout the paper.
\end{EN}

\begin{physbox}{为什么光谱合成天然适合 GPU}
辐射转移只在同一波长上耦合不同的大气深度（光子在某一频率上穿过各层被吸收和再发射），但一个波长的计算几乎不影响另一个波长。于是“一百万个波长”就是一百万份互相独立的短计算——这正是 GPU 最擅长的大规模并行模式。原 Kurucz 程序诞生于内存极小的年代，被迫逐深度写盘、逐波长串行求解；本文只是换了执行方式，让同一个物理方程以大批量数组运算的形式在 GPU 上完成。物理内容一行未改，验证残差在 $10^{-3}$ 量级以下。
\end{physbox}

\begin{notebox}{表~\ref{tab:atmosphere_symbols} 说明}
下表为原文 Table 1 的中英对照精简版，按用途分组列出全文符号；定义栏先中文后英文（英文有所节略）。
\end{notebox}

\begingroup
\footnotesize
\begin{longtable}{@{}p{0.16\textwidth}p{0.28\textwidth}p{0.48\textwidth}@{}}
\caption{全文符号与定义（中英对照）。Symbols and definitions used throughout the paper.}
\label{tab:atmosphere_symbols}\\
\toprule
\textbf{角色 Role} & \textbf{符号 Symbol} & \textbf{定义 Definition} \\
\midrule
\endfirsthead
\toprule
\textbf{角色 Role} & \textbf{符号 Symbol} & \textbf{定义 Definition} \\
\midrule
\endhead
恒星标签 & $T_{\rm eff}$, $\log g$, $g$, $\xi$, $[\mathrm{M}/\mathrm{H}]$, $[\alpha/\mathrm{M}]$, $[\mathrm{Fe}/\mathrm{H}]$, $\{[\mathrm{X}/\mathrm{H}]\}$, $\{[\mathrm{X}/\mathrm{Fe}]\}$, $A(\mathrm{X})$ & 有效温度、表面重力标签、物理加速度 $g=10^{\log g}$（cgs）、微观湍流与丰度。$[\mathrm{M}/\mathrm{H}]$ 整体移动所有金属，$[\alpha/\mathrm{M}]$ 额外叠加于 O、Ne、Mg、Si、S、Ca、Ti。$[\mathrm{X}/\mathrm{H}]\equiv\log_{10}(n_{\rm X}/n_{\rm H})-\log_{10}(n_{\rm X}/n_{\rm H})_\odot$；$[\mathrm{X}/\mathrm{Fe}]=[\mathrm{X}/\mathrm{H}]-[\mathrm{Fe}/\mathrm{H}]$；$A(\mathrm{X})=\log_{10}(n_{\rm X}/n_{\rm H})+12$。显式元素坐标优先于整体或分组缩放。Effective temperature, gravity, microturbulence and abundances; an explicit element coordinate overrides bulk scaling. \\
深度坐标 & $m$, $\tau_{\rm R}$ & 柱质量向内增加；罗斯兰光学深度满足 $\mathrm{d}\tau_{\rm R}=\kappa_{\rm R}\,\mathrm{d}m$。Column mass and Rosseland optical depth. \\
大气状态 & $T$, $\rho$, $P_{\rm gas}$, $n_e$, $\kappa_{\rm R}$, $g_{\rm rad}$, $\xi$ & 各深度层的温度、质量密度、气体压强、电子密度、罗斯兰不透明度与辐射加速度，外加与深度无关的微观湍流速度。 \\
压强项 & $P_{\rm rad}$, $P_0$ & 辐射压强，以及无湍流压强分支中 $P_{\rm gas}+P_{\rm rad}=gm+P_0$ 的表面常数。 \\
粒子布居 & $n_{s,r}$, $U_{s,r}$, $n_{\rm mol}$, $n_{\rm pert}$, $\mathcal{I}_{s,r}$, $\Delta\mathcal{I}_{s,r}$ & 物种 $s$ 第 $r$ 电离级的数密度与配分函数、分子数密度、中性碰撞代理量、电离能及其随密度变化的下降修正。 \\
能量平衡 & $F_{\rm rad}$, $F_{\rm conv}$, $\sigma_{\rm SB}$, $\epsilon_F$ & 辐射流量与对流流量满足 $F_{\rm rad}+F_{\rm conv}=\sigma_{\rm SB}T_{\rm eff}^4$；$\epsilon_F$ 为相对残差。 \\
迭代 & $n$, $j$, $\Delta T^{(n)}$, $\epsilon_T$ & 迭代序号、深度序号、温度改正量，以及深层最大相对温度变化。 \\
不动点 & $\mathbf{x}^{(n)}$, $\mathbf{p}^{(n)}$, $\mathcal{G}$, $\mathbf{J}_{\mathcal{G}}$, $\mathcal{R}$, $\mathbf{x}_\star$, $\mathbf{e}^{(n)}$ & 数值状态及其六廓线投影、一个大气循环、局部雅可比与残差、收敛不动点与迭代误差。 \\
隐变量初始化器 & $\mathbf{u}$, $\overline{\mathbf u}$, $\mathbf{s}_u$, $\widehat{\mathbf z}$, $\widehat{\mathbf c}$, $\widehat{\mathbf u}$, $\mathbf{B}_K$, $K$, $\mathcal{D}$, $\Delta m_j$, $T_{\rm grey}$, $s_g$ & 变换后的廓线及其训练均值与标度、标准化网络预测、反标准化 PCA 系数、解码廓线、基与维度、解码器，以及质量、温度与加速度变换。 \\
训练损失 & $\mathcal{L}$, $\mathcal{L}_{\rm prof}$, $\mathcal{L}_{\nabla}$, $\mathcal{L}_{\tau}$, $\mathcal{L}_{\rm hse}$, $\lambda$ & 初始化器总目标及其解码廓线、深度导数、光学深度与流体静力学各项及其权重。 \\
约化求解器 & $T(\tau_{\rm R})$, $m(\tau_{\rm R})$, $\mathcal{R}_T$, $\mathcal{R}_m$ & 独立的温度与柱质量廓线及其能量平衡、光学深度残差。 \\
常数 & $c$, $h$, $k_{\rm B}$, $e$, $m_e$ & 光速、普朗克与玻尔兹曼常数、元电荷、电子质量。 \\
原子谱线 & $l$, $\nu_l$, $(gf)_l$, $E_l$, $m_l$ & 跃迁序号、频率、统计权重振子强度、下能级能量、吸收体质量。 \\
致宽 & $\gamma_{\rm rad}$, $\gamma_{\rm S}$, $\gamma_{\rm vdW}$ & 辐射阻尼率，以及每个扰动粒子的斯塔克、范德瓦尔斯系数（分别乘电子与中性碰撞密度）。 \\
谱线定标 & $q(l)$, $\beta$, $\delta_{q,gf}$, $\delta_{q,\beta}$ & 谱线分量 $l$ 所属的改正组、阻尼族序号、对振子强度与阻尼的有界改正；撇号表示定标后的值。 \\
采样 & $\nu$, $\lambda_{i_\lambda}$, $\lambda_0$, $\lambda_{\min}$, $\lambda_{\max}$, $R_{\rm grid}$, $N_\lambda$ & 频率、波长坐标、首个采样与端点。合成网格满足 $\lambda_{i_\lambda+1}/\lambda_{i_\lambda}=1+R_{\rm grid}^{-1}$，共 $N_\lambda$ 个采样。 \\
轮廓 & $\Delta v_{\rm D}$, $\Delta\nu_{\rm D}$, $a_l$, $\phi_l$ & 多普勒速度宽度与频率宽度、阻尼比，以及按 $\int\phi_l(\nu)\,\mathrm{d}\nu=1$ 归一的谱线轮廓。 \\
不透明度 & $\kappa_{\nu,{\rm c}}$, $\sigma_{\nu,{\rm c}}$, $\kappa_{\nu,l}$, $\chi_\nu$ & 连续吸收、连续散射与谱线吸收合成 $\chi_\nu=\kappa_{\nu,{\rm c}}+\sigma_{\nu,{\rm c}}+\sum_l\kappa_{\nu,l}$。 \\
谱线筛选 & $b$, $\nu_b$, $\eta_{\rm sel}$, $C_{jb}$, $Q_{srb}$, $C_0$, $R_l$, $T_*$ & 频率 bin 序号与参考频率、筛选比例、连续谱阈值、布居/连续谱代理最大值、固定常数与中心线强代理；$T_*$ 为最深层温度。 \\
辐射转移 & $\tau_\nu$, $\epsilon_\nu$, $\mu$, $B_\nu$, $J_\nu$, $S_\nu$, $I_\nu$ & 光学深度、热吸收占比、外向方向余弦、普朗克函数、平均强度、源函数与比强度。 \\
输出 & $H_\nu$, $H_\lambda$, $F_\nu$, $F_\lambda$, $F_{\lambda,{\rm c}}$, $f_{\rm norm}$ & $H$ 为一阶角矩，$F=4\pi H$。表面流量即光谱，$F_{\lambda,{\rm c}}$ 为其连续谱，$f_{\rm norm}$ 为二者之比。 \\
仪器 & $\mathcal{K}$ & 从内禀总流量与连续流量到观测网格的线扩展函数线性算子。 \\
光谱拟合 & $\boldsymbol{\theta}$, $y_p$, $w_p$, $f_p$, $f_{\rm fast}$, $f_{\rm phys}$, $\Delta f$, $Q$, $\delta Q_{\min}$ & 拟合标签、观测归一化流量、逆方差权重、像素 $p$ 处模型流量。$f_{\rm fast}$/$f_{\rm phys}$ 为大气收敛前/后的光谱，$Q$ 为每有效像素的平均卡方。 \\
巡天冗余参数 & $\boldsymbol{\psi}$, $v_r$, $v_b$, $\mathcal{S}$, $\mathcal{B}$ & 联合拟合状态、残余速度、有效高斯致宽、多普勒算子与致宽算子。 \\
\bottomrule
\end{longtable}
\endgroup

\subsection{物理计算 The physical calculation}

我们在对数波长上均匀采样合成网格：
\begin{equation}
\lambda_{i_\lambda} =
\lambda_0\left(1+R_{\rm grid}^{-1}\right)^{i_\lambda} .
\end{equation}

\noindent 参数 $R_{\rm grid}$ 设定的是模型采样密度，而非仪器分辨本领。仪器模型通过 $\mathcal{K}$ 施加线扩展函数，一个分辨率元内可以容纳多个模型采样点。采样数约为
\begin{equation}
N_\lambda \simeq
\frac{\ln(\lambda_{\max}/\lambda_{\min})}
     {\ln(1+R_{\rm grid}^{-1})}
\simeq R_{\rm grid}\ln\left(\frac{\lambda_{\max}}{\lambda_{\min}}\right).
\end{equation}

\begin{EN}
We sample the synthesis grid uniformly in logarithmic wavelength,
\begin{equation}
\lambda_{i_\lambda} =
\lambda_0\left(1+R_{\rm grid}^{-1}\right)^{i_\lambda} .
\end{equation}

\noindent The parameter $R_{\rm grid}$ sets the model sampling, not the instrumental
resolving power. The instrument model applies the line spread function through
$\mathcal{K}$, and several model samples can lie within one resolution element.
The number of samples is approximately
\begin{equation}
N_\lambda \simeq
\frac{\ln(\lambda_{\max}/\lambda_{\min})}
     {\ln(1+R_{\rm grid}^{-1})}
\simeq R_{\rm grid}\ln\left(\frac{\lambda_{\max}}{\lambda_{\min}}\right).
\end{equation}
\end{EN}

在每个波长处，质量消光系数是连续吸收、连续散射与谱线吸收之和：
\begin{equation}
\chi_\nu = \kappa_{\nu,{\rm c}} + \sigma_{\nu,{\rm c}}
          + \sum_l \kappa_{\nu,l} .
\end{equation}

\noindent 连续项包括来自氢、氦和金属的主要束缚—自由与自由—自由过程 \citep{Mihalas1978,Gray2005}。在冷星的光学与近红外连续谱中，H$^-$ 占据主导地位 \citep{John1988}。瑞利散射与电子散射单独处理，因为它们的来源不是纯热辐射。

\begin{EN}
At each wavelength the mass extinction coefficient is the sum of continuous
absorption, continuous scattering, and line absorption,
\begin{equation}
\chi_\nu = \kappa_{\nu,{\rm c}} + \sigma_{\nu,{\rm c}}
          + \sum_l \kappa_{\nu,l} .
\end{equation}

\noindent The continuous terms include the principal bound-free and free-free
processes from hydrogen, helium, and metals \citep{Mihalas1978,Gray2005}.
H$^-$ dominates much of the optical and near-infrared continuum in cool stars
\citep{John1988}. Rayleigh and electron scattering remain separate because
their source is not purely thermal.
\end{EN}

\begin{physbox}{不透明度与 H$^-$ 的故事}
不透明度 $\chi_\nu$ 描述气体在某频率上“挡光”的能力，分三部分：连续吸收（光致电离、自由—自由跃迁）、连续散射（瑞利散射、电子散射）、以及所有谱线吸收之和。其中最有趣的是 H$^-$：一个中性氢原子临时“挂靠”一个额外电子形成负离子，它可以吸收很宽频率范围的光子。在冷星光球里氢几乎不电离，自由电子稀缺且主要来自金属——这解释了前文“金属丰度改变连续谱”的通道：金属 $\rightarrow$ 电子 $\rightarrow$ H$^-$ $\rightarrow$ 连续谱。H$^-$ 是 1938 年才被 Wildt 提出、并被证认为太阳光学连续谱的主要来源，是恒星大气物理中“看似次要、实则主导”的经典案例。
\end{physbox}

对于具有频率归一轮廓的标准原子跃迁，LTE 下单位质量的谱线不透明度为 \citep{Gray2005}
\begin{equation}
\begin{aligned}
\kappa_{\nu,l}={}&
\frac{\pi e^2}{m_e c}\,
\frac{n_{s,r}}{\rho U_{s,r}}\,(gf)_l
\exp\left(-\frac{E_l}{k_{\rm B}T}\right) \\
&{}\times
\left[1-\exp\left(-\frac{h\nu}{k_{\rm B}T}\right)\right]\phi_l(\nu), \\
&\text{with}\quad \int\phi_l(\nu)\,\mathrm{d}\nu=1 .
\end{aligned}
\end{equation}

\noindent 其中 $n_{s,r}/U_{s,r}$ 是物种 $s$ 第 $r$ 电离级经配分函数归一的数密度，$(gf)_l$ 包含下能级统计权重。轮廓宽度与阻尼尺度可概括为
\begin{equation}
\begin{aligned}
\Delta v_{\rm D}^{2} &=
\frac{2k_{\rm B}T}{m_l}+\xi^2, \\
\Delta\nu_{\rm D} &= \frac{\nu_l}{c}\Delta v_{\rm D}, \\
a_l &=
\frac{\gamma_{\rm rad}+\gamma_{\rm S}n_e+
\gamma_{\rm vdW}n_{\rm pert}}{4\pi\Delta\nu_{\rm D}} .
\end{aligned}
\end{equation}

\noindent 这里 $n_{\rm pert}$ 是范德瓦尔斯处方使用的有效中性碰撞密度代理。得到的福格特轮廓决定中心不透明度在波长上的分布。在布居固定时，振子强度缩放频率积分的不透明度，而阻尼把不透明度从线心重新分配到翼部 \citep{Barklem2000,Gray2005}。温度、电子密度、丰度和电离状态既改变吸收体的数量，也改变它们所处的致宽环境。这些直接联系使得恒星标签与原子数据可以在光谱空间中被优化。

\begin{EN}
For a standard atomic transition with a frequency normalized profile, the LTE
line opacity per unit mass is \citep{Gray2005}
\begin{equation}
\begin{aligned}
\kappa_{\nu,l}={}&
\frac{\pi e^2}{m_e c}\,
\frac{n_{s,r}}{\rho U_{s,r}}\,(gf)_l
\exp\left(-\frac{E_l}{k_{\rm B}T}\right) \\
&{}\times
\left[1-\exp\left(-\frac{h\nu}{k_{\rm B}T}\right)\right]\phi_l(\nu), \\
&\text{with}\quad \int\phi_l(\nu)\,\mathrm{d}\nu=1 .
\end{aligned}
\end{equation}

\noindent Here $n_{s,r}/U_{s,r}$ is the partition-normalized number density of
ion stage $r$ of species $s$, and $(gf)_l$ includes the lower-level statistical
weight. The profile width and
damping scale can be summarized by
\begin{equation}
\begin{aligned}
\Delta v_{\rm D}^{2} &=
\frac{2k_{\rm B}T}{m_l}+\xi^2, \\
\Delta\nu_{\rm D} &= \frac{\nu_l}{c}\Delta v_{\rm D}, \\
a_l &=
\frac{\gamma_{\rm rad}+\gamma_{\rm S}n_e+
\gamma_{\rm vdW}n_{\rm pert}}{4\pi\Delta\nu_{\rm D}} .
\end{aligned}
\end{equation}

\noindent Here $n_{\rm pert}$ is the effective neutral-collision density proxy
used by the van der Waals prescription.
The resulting Voigt profile sets how the central opacity is distributed in
wavelength. At fixed population, the oscillator strength scales the
frequency-integrated opacity, while damping redistributes opacity from the
core into the wings \citep{Barklem2000,Gray2005}. Temperature, electron density, abundance,
and ionization change both the number of absorbers and their broadening
environment. These direct connections permit stellar labels and atomic data
to be optimized in spectrum space.
\end{EN}

\begin{physbox}{一条谱线是怎么形成的：振子强度、致宽与福格特轮廓}
上式是全文最核心的物理公式，拆解如下：
\begin{itemize}
\item \textbf{吸收体数量}：$\exp(-E_l/k_{\rm B}T)$ 是玻尔兹曼因子——下能级能量 $E_l$ 越高，能处于该能级的原子越少；配分函数 $U_{s,r}$ 则归一化各能级的布居。温度既改变激发也改变电离，所以同一条谱线的强度随恒星类型剧烈变化，这正是光谱型分类的物理基础。
\item \textbf{振子强度 $(gf)$}：量子力学给出的跃迁“固有强度”，缩放整条谱线的积分吸收。它是本文第 6 节定标的主角——许多谱线的 $gf$ 值理论或实验室误差不小，要靠太阳等基准星反推。
\item \textbf{受激辐射修正}：$[1-\exp(-h\nu/k_{\rm B}T)]$ 是 LTE 下扣掉受激发射的因子。
\item \textbf{轮廓 $\phi_l(\nu)$}：多普勒致宽（热运动 $\propto\sqrt{T/m}$ 加微观湍流 $\xi$）给出高斯线心；阻尼（自然宽度、电子碰撞斯塔克致宽、中性原子碰撞范德瓦尔斯致宽）给出洛伦兹翼部。高斯与洛伦兹的卷积就是\textbf{福格特轮廓}。阻尼参数决定强线翼部的形状，是测量表面重力的重要探针（巨星气压低、阻尼翼弱，矮星反之）。
\end{itemize}
\end{physbox}

标准福格特轮廓并不足以描述所有跃迁。我们实现了氢的精细结构、斯塔克轮廓 \citep{VidalCooperSmith1973}，以及通过占据概率物理把高序谱线并入连续谱的处理 \citep{HummerMihalas1988}。独立的 Kurucz 分支处理 He~I、$^3$He、He~II 以及不对称的 Shore--Fano 自电离轮廓 \citep{Fano1961,Shore1967}。这些专门轮廓遵循 \citet{KuruczAvrett1981}。分子谱线不透明度使用分子平衡给出的布居 \citep{Tsuji1973}，分子质量随谱线数据存储并用于设定其多普勒宽度。

\begin{EN}
The standard Voigt profile is not sufficient for every transition. We
implement hydrogen fine structure, Stark profiles
\citep{VidalCooperSmith1973}, and the merging of high series members into the
continuum through occupation-probability physics \citep{HummerMihalas1988}.
Separate Kurucz branches treat He~I, $^3$He, He~II, and asymmetric
Shore--Fano autoionizing profiles \citep{Fano1961,Shore1967}. The specialized
profiles follow \citet{KuruczAvrett1981}.
Molecular line opacity uses populations from molecular equilibrium
\citep{Tsuji1973} and molecular masses stored with the line data
to set their Doppler widths.
\end{EN}

消光通过 $\mathrm{d}\tau_\nu=\chi_\nu\,\mathrm{d}m$ 定义光学深度，其中 $m$ 为柱质量。LTE 下热辐射源是普朗克函数，谱线不透明度按热吸收处理；连续散射的源则依赖平均强度 \citep{Mihalas1978,Sobeck2011}。因此源函数可写为
\begin{equation}
\epsilon_\nu =
\frac{\kappa_{\nu,{\rm c}}+\sum_l\kappa_{\nu,l}}{\chi_\nu},
\qquad
S_\nu = \epsilon_\nu B_\nu(T)
      + \left(1-\epsilon_\nu\right)J_\nu .
\end{equation}

对每个波长，把热源与散射占比映射到光学深度之后，我们求解
\begin{equation}
\mu\frac{\mathrm{d}I_\nu}{\mathrm{d}\tau_\nu}=I_\nu-S_\nu .
\end{equation}

\noindent 转移计算对散射贡献做迭代。由于各波长保持独立，整个波长网格可以并行求解。取 $\mu$ 相对向外法线计量，出射的爱丁顿流量经下式换算：
\begin{equation}
F_\lambda
=F_\nu\left|\frac{\mathrm{d}\nu}{\mathrm{d}\lambda}\right|
=4\pi H_\nu\frac{c}{\lambda^2}
=4\pi H_\lambda .
\end{equation}
\noindent 报告的绝对流量为每纳米的 $F_\lambda$，其中 $c$ 与 $\lambda$ 以一致的纳米单位取值。我们对每个波长携带两条转移分支：一条包含谱线与连续不透明度，另一条只含连续谱。二者一起求解，给出总流量 $F_\lambda$ 及其物理连续谱 $F_{\lambda,{\rm c}}$。存在仪器算子时，先对两者分别施加仪器算子再取比值：
\begin{equation}
f_{\rm norm} =
\frac{\mathcal{K}F_\lambda}
     {\mathcal{K}F_{\lambda,{\rm c}}} .
\end{equation}

\noindent 只要连续谱在一个分辨率元内有变化，这个顺序就很重要。它还让优化器可以直接比较模型与观测像素，而不必在合成与卷积之间把密集采样的光谱搬回 CPU 内存。以上方程定义了物理计算本身；实现层面改变的只是这些运算的组织与执行方式。

\begin{EN}
The extinction defines optical depth through
$\mathrm{d}\tau_\nu=\chi_\nu\,\mathrm{d}m$, where $m$ is column mass. In LTE the thermal source
is the Planck function and line opacity is treated as thermal absorption. The
continuum scattering source depends on the mean intensity
\citep{Mihalas1978,Sobeck2011}. The source function
can therefore be written as
\begin{equation}
\epsilon_\nu =
\frac{\kappa_{\nu,{\rm c}}+\sum_l\kappa_{\nu,l}}{\chi_\nu},
\qquad
S_\nu = \epsilon_\nu B_\nu(T)
      + \left(1-\epsilon_\nu\right)J_\nu .
\end{equation}

For each wavelength, after mapping the thermal source and scattering fraction
to optical depth, we solve
\begin{equation}
\mu\frac{\mathrm{d}I_\nu}{\mathrm{d}\tau_\nu}=I_\nu-S_\nu .
\end{equation}

\noindent The transfer calculation iterates the scattering contribution. Since
wavelengths remain independent, the full wavelength grid can be solved in
parallel. With $\mu$ measured from the outward surface normal, the emergent
Eddington flux is converted through
\begin{equation}
F_\lambda
=F_\nu\left|\frac{\mathrm{d}\nu}{\mathrm{d}\lambda}\right|
=4\pi H_\nu\frac{c}{\lambda^2}
=4\pi H_\lambda .
\end{equation}
\noindent The reported absolute flux is $F_\lambda$ per nanometer, with $c$
and $\lambda$ evaluated in consistent nanometer units.
We carry two transfer branches for each wavelength. One
contains line and continuum opacity. The other contains only the continuum.
We solve them together to produce the total flux $F_\lambda$ and its physical
continuum $F_{\lambda,{\rm c}}$. When an instrument operator is present, we
apply it to both before taking their ratio,
\begin{equation}
f_{\rm norm} =
\frac{\mathcal{K}F_\lambda}
     {\mathcal{K}F_{\lambda,{\rm c}}} .
\end{equation}

\noindent This order matters whenever the continuum varies across a resolution element.
It also lets an optimizer compare the model directly with observed pixels
without moving a densely sampled spectrum back to CPU memory between synthesis
and convolution. These equations define the physical calculation. The
implementation changes how the same operations are arranged and executed.
\end{EN}

\begin{physbox}{源函数、辐射转移方程与 LTE}
辐射转移方程 $\mu\,\mathrm{d}I_\nu/\mathrm{d}\tau_\nu=I_\nu-S_\nu$ 的含义很直观：一束光穿过气体时，强度的净变化 = 被吸收的（$-I_\nu$）加上新发射的（$+S_\nu$）。全部物理都装在\textbf{源函数} $S_\nu$ 里。在 LTE（局部热动平衡）近似下，热吸收对应的辐射由当地温度的普朗克函数 $B_\nu(T)$ 给出——即假设气体局部就像一个温度为 $T$ 的小黑体，碰撞足够频繁，使能级布居由温度唯一决定，无需逐个求解辐射与能级的耦合（那是 NLTE 计算的昂贵所在）。散射部分则不同：散射光子只是“换了个方向”，其辐射强度取决于各方向的平均强度 $J_\nu$，所以要做迭代。我们看到的光谱，就是各深度辐射沿视线方向累积出来的结果：光学深度 $\tau_\nu\sim1$ 的层次贡献最大——谱线中心不透明度大，其 $\tau=1$ 面更高、更冷，所以谱线表现为连续谱背景上的吸收凹陷。
\end{physbox}

\begin{figure}[t]
\centering
\includegraphics[width=0.98\textwidth]{figure_01.pdf}
\caption{与原 Kurucz 计算在 300--1000 nm、$R_{\rm grid}=300{,}000$ 下的一致性。各行依次为热矮星、太阳、红巨星与 K 矮星；各列依次为完整归一化光谱、一个代表性特征，以及 $\Delta f_{\rm norm}$ 残差。黑色与橙色比较两套完整计算；残差面板中蓝色固定大气以单独检验合成，橙色包含在相同标签下独立收敛的大气。残差乘以 $10^3$。自上而下，仅合成与完整计算的 99 百分位绝对残差（单位 $10^{-3}$）分别为 $(2.0,0.7)$、$(2.8,2.2)$、$(5.8,5.8)$ 和 $(3.0,2.9)$。（Agreement with the original Kurucz calculation over 300--1000 nm.）}
\label{fig:synthesis_overview}
\end{figure}

\subsection{从分阶段文件到波长并行执行 From staged files to wavelength-parallel execution}

\begin{figure}[t]
\centering
\includegraphics[width=0.98\textwidth]{figure_02.pdf}
\caption{$R_{\rm grid}=300{,}000$ 下每条光谱的合成运行时间。左图比较四种恒星类型在原始串行 Kurucz 程序与 \paynezero（单张 V100、A100 或 H100 GPU）上的 300--1000 nm 计算。中上、右上面板分解太阳与 K 矮星的 GPU 耗时，下面板给出 1500--1700 nm 太阳光谱的总耗时。原子与分子谱线不透明度主导 GPU 运行时间。（Synthesis runtime per spectrum.）}
\label{fig:synthesis_runtime}
\end{figure}

\begin{figure}[t]
\centering
\includegraphics[width=0.8\textwidth]{figure_03.pdf}
\caption{合成耗时随波长跨度的缩放（$R_{\rm grid}=300{,}000$）。所有窗口从 300 nm 起，止于 400、500、700 或 1000 nm。彩色曲线为 H100 上四类恒星；虚线与阴影带为原 Kurucz 耗时的中位数与范围。更平缓的 H100 缩放使加速比随波长区间与保留跃迁数增大。（Synthesis scaling with wavelength span.）}
\label{fig:synthesis_scaling}
\end{figure}

原始的 Kurucz 计算以一系列 Fortran 程序串行运行。谱线读取程序先准备所请求的波长区间；不透明度阶段随后沿大气深度循环，每一层都重绕谱线记录、累积谱线不透明度并写出按深度排序的记录；直接存取文件把这些记录转置；转移阶段读取转置后的不透明度，逐波长推进，在每个波长点求解连续谱与总量的转移问题。这种组织方式与当年程序编写时的内存条件非常匹配，但也产生了中间输入输出，并让最大的独立维度基本保持串行 \citep{Kurucz2005}。

\begin{EN}
The original Kurucz calculation runs as a sequence of Fortran programs. Line
readers first prepare a requested wavelength interval. The opacity stage then
loops over atmosphere depth, rewinds the line records for each layer,
accumulates line opacity, and writes depth-ordered records. Direct-access files
transpose these records. The transfer stage reads the transposed opacity and
advances through wavelength one point at a time. At every point it solves the
continuum and total transfer problems. This organization was well
matched to the memory available when the programs were written. It also
creates intermediate input and output and leaves the largest independent
dimension largely serial \citep{Kurucz2005}.
\end{EN}

并行执行的机会来自物理依赖关系。辐射转移在一个波长上耦合各层大气，但并不把一个波长耦合到下一个波长。因此自然的工作单元是一大组相互独立的波长列，每列只做一次短深度计算。历史上按深度排序的记录和文件转置掩盖了这一结构；把完整的深度—波长数组保留在内存中，则直接暴露了它。

\begin{EN}
The opportunity for parallel execution follows from the physical dependencies.
Radiative transfer couples the atmosphere depths at one wavelength, but it does
not couple one wavelength to the next. The natural unit of work is therefore a
large set of independent wavelength columns, each with a short calculation
through depth. The historical depth-ordered records and file transpose obscure
this structure. Keeping the complete depth and wavelength array in memory
exposes it directly.
\end{EN}

当一个编译算子同时作用于大量数组元素时，GPU 才能发挥效力，这样的算子常被称为 kernel。在 \paynezero\ 中，我们传给每个 kernel 的是一大组波长或活跃谱线，而非单个波长。物理方程保持不变，但反复的小调用被合并成编译后的数组运算。连续不透明度覆盖整个深度—波长数组，活跃谱线并行地叠加其轮廓，我们对所有波长沿短深度轴积分光学深度。总量与连续谱转移共享同一批次。所有数组和固定的转移几何一直留在 GPU 上，直到 kernel 返回最终流量。

\begin{EN}
A GPU is effective when one compiled operation acts on many array
elements at once. Such an operation is often called a kernel. In \paynezero,
we pass each kernel a large group of wavelengths or active lines rather than
one wavelength. We leave the physical equations unchanged, but repeated
small calls become compiled array operations. Continuum opacity spans the full
depth and wavelength array, active lines add their profiles in parallel, and
we integrate optical depth along the short depth axis for all wavelengths.
Total and continuum transfer share the same batch. We keep the arrays and fixed
transfer geometry on the GPU until the kernel returns the final flux.
\end{EN}

因此，占主导地位的窄谱线与波长计算避免了对单个轮廓或波长采样的 Python 循环。专门的宽线与自电离轮廓以及有界的目录块仍保留少量 Python 层面的调度循环。PyTorch 把大数组运算发往 CUDA 或 CPU kernel，而分子目录由一个缓存的 Numba kernel 准备。测试过的 Apple Metal 路径使合成可以在 Apple silicon 本地运行。所报告的耗时与吞吐建议均基于专用 NVIDIA V100、A100、H100 加速卡。由于 Kurucz 程序本身也是编译代码，速度差异来自并发、更大的 kernel 和常驻输入，而非仅仅来自编译。表~\ref{tab:synthesis_speed} 总结了这些改变如何替代串行循环与中间文件。在反复拟合过程中，固定的窗口、目录与转移几何保持常驻，而依赖大气的不透明度与流量对每个标签向量重新计算。不缓存、不模拟任何合成流量。

\begin{EN}
The dominant narrow-line and wavelength calculations therefore avoid Python
loops over individual profiles or wavelength samples. Specialized broad and
autoionizing profiles and bounded catalog chunks retain small Python-level
dispatch loops. PyTorch sends their large array operations to CUDA or CPU
kernels, while a cached Numba kernel prepares the molecular catalog. The tested Apple Metal path makes synthesis available
locally on Apple silicon. The reported timings and throughput
recommendations use dedicated NVIDIA V100, A100, and H100 accelerators. Since the
Kurucz programs are also compiled, the speed difference comes from concurrency,
larger kernels, and resident inputs rather than compilation alone.
Table~\ref{tab:synthesis_speed} summarizes how these changes replace serial
loops and intermediate files. During repeated fits, the fixed window, catalog,
and transfer geometry remain resident, while atmosphere-dependent opacity and
flux are recomputed for every label vector. No synthetic flux is cached or
emulated.
\end{EN}

\begin{notebox}{表~\ref{tab:synthesis_speed} 说明（原文 Table 2，内容保留英文）}
该表逐阶段对比原 Kurucz 与 \paynezero\ 的执行方式：窗口准备（Numba 编译 + 进程内缓存）、工作布局（PyTorch 批量波长）、谱线不透明度（有界块 + 中心强度测试剔除弱“深度—谱线”对）、内存与 I/O（深度—波长数组常驻设备，取消磁盘中转）、辐射转移（所有波长与两条转移分支一次张量计算）。
\end{notebox}

\begin{table*}[t]
\centering
\caption{Execution of the same synthesis stages in the original Kurucz
workflow and \paynezero. The final column summarizes how resident data and
wavelength-parallel GPU kernels change the computational cost.}
\label{tab:synthesis_speed}
\begingroup
\setlength{\tabcolsep}{4pt}
\renewcommand{\arraystretch}{1.18}
\footnotesize
\begin{tabular}{@{}llll@{}}
\toprule
\tcell{0.12\textwidth}{\textbf{Stage}} &
\tcell{0.25\textwidth}{\textbf{Original Kurucz execution}} &
\tcell{0.28\textwidth}{\textbf{\paynezero\ execution}} &
\tcell{0.25\textwidth}{\textbf{Effect}} \\
\midrule
\tcell{0.12\textwidth}{\textbf{Window preparation}} &
\tcell{0.25\textwidth}{Line readers prepare wavelength-specific sequential
records for the staged programs} &
\tcell{0.28\textwidth}{Numba compiles molecular source records into portable
window-bound arrays. An in-process cache retains the grid, catalogs, tables,
and fixed tensors for one device, dtype, and chunk configuration} &
\tcell{0.25\textwidth}{The retained measurements of repeated calls amortize this
preparation without making the disk products device specific} \\
\tcell{0.12\textwidth}{\textbf{Work layout}} &
\tcell{0.25\textwidth}{Compiled Fortran routines retain serial outer loops over
depth records and wavelengths} &
\tcell{0.28\textwidth}{PyTorch batches independent wavelengths and the dominant
narrow-line and depth contributions. The retained V100, A100, and H100
measurements use CUDA. The same path also supports Metal and CPU execution} &
\tcell{0.25\textwidth}{The bulk work advances concurrently, while Python dispatch is
limited to bounded chunks and specialized profiles} \\
\tcell{0.12\textwidth}{\textbf{Line opacity}} &
\tcell{0.25\textwidth}{Each depth rereads atomic and molecular line records and
advances through profiles in record order} &
\tcell{0.28\textwidth}{Atomic and molecular catalogs are processed in bounded
line chunks. A continuum-referenced center test rejects weak depth--line pairs,
and survivors deposit only their required wing spans} &
\tcell{0.25\textwidth}{Cost and temporary memory follow retained-line scans and
active profile deposits, not the full line-list--wavelength product} \\
\tcell{0.12\textwidth}{\textbf{Memory and I/O}} &
\tcell{0.25\textwidth}{Depth-ordered opacity is written and transposed through
intermediate files} &
\tcell{0.28\textwidth}{Depth--wavelength opacity, source, and transfer arrays
stay on the device between stages, while line and active-pair chunks bound temporary
memory} &
\tcell{0.25\textwidth}{Disk opacity staging and transfers of full arrays through CPU memory are
removed without constructing a depth--line--wavelength array} \\
\tcell{0.12\textwidth}{\textbf{Transfer}} &
\tcell{0.25\textwidth}{Wavelengths advance serially, with separate total and
continuum calls} &
\tcell{0.28\textwidth}{All wavelengths and the stacked total and continuum
branches share one tensor calculation} &
\tcell{0.25\textwidth}{The largest independent dimension is evaluated
together} \\
\addlinespace[2pt]
\bottomrule
\end{tabular}
\endgroup
\end{table*}

\subsection{精度与性能 Accuracy and performance}

我们在 $R_{\rm grid}=300{,}000$ 的对数网格上检验 300--1000 nm 的最终光谱。图~\ref{fig:synthesis_overview} 做了两组比较：共享大气与谱线目录用于单独检验合成精度；在相同恒星标签下各自独立计算则把大气与合成一起检验。两者都以原始 Fortran Kurucz 链条为基准。\footnote{\url{https://github.com/tingyuansen/kurucz}}

\begin{EN}
We test the final spectrum over 300--1000~nm on a logarithmic grid with
$R_{\rm grid}=300{,}000$. Figure~\ref{fig:synthesis_overview} makes two
comparisons. A shared atmosphere and line catalog isolate synthesis accuracy.
Independent calculations at the same stellar labels test atmosphere and
synthesis together. Both use the original Fortran Kurucz chain as the
reference.\footnote{\url{https://github.com/tingyuansen/kurucz}}
\end{EN}

跨四种恒星类型，归一化流量残差的合并均方根在仅合成时为 $9.1\times10^{-4}$，完整计算为 $8.6\times10^{-4}$。对热矮星、太阳、红巨星和 K 矮星，99 百分位绝对残差仅合成时分别为 $(2.0,2.8,5.8,3.0)\times10^{-3}$，完整计算时为 $(0.7,2.2,5.8,2.9)\times10^{-3}$。这些残差反映的是同一物理计算的两种实现之间的微小数值差异。

\begin{EN}
Across the four stellar types, the pooled rms normalized-flux residual is
$9.1\times10^{-4}$ for synthesis alone and $8.6\times10^{-4}$ for the complete
calculation. For the hot dwarf, Sun, red giant, and K dwarf, respectively, the
99th-percentile absolute residuals are $(2.0,2.8,5.8,3.0)\times10^{-3}$ for
synthesis alone and $(0.7,2.2,5.8,2.9)\times10^{-3}$ for the complete
calculation. These residuals reflect small numerical differences between two
implementations of the same physical calculation.
\end{EN}

图~\ref{fig:synthesis_runtime} 测量的是一次性源窗口准备之后的重复前向计算。每次计时都包含连续不透明度、原子与分子谱线、总量与连续谱转移以及返回的光谱。在 $R_{\rm grid}=300{,}000$ 的 300--1000 nm 网格上，太阳光谱的中位耗时约为：原 Kurucz 程序 763 秒、单张 V100 为 32 秒、A100 为 29 秒、H100 为 14 秒，对应加速比分别为 24、26 与 54 倍。

\begin{EN}
Figure~\ref{fig:synthesis_runtime} measures repeated forward calculations after
one-time source-window preparation. Each wall includes continuum opacity,
atomic and molecular lines, total and continuum transfer, and the returned
spectrum. On the 300--1000 nm grid at $R_{\rm grid}=300{,}000$, the solar
medians are about 763~s for the original Kurucz programs, 32~s for one V100,
29~s for one A100, and 14~s for one H100. These correspond to speedups of
$24$, $26$, and $54$ on the V100, A100, and H100, respectively.
\end{EN}

对同一太阳大气、较窄的 APOGEE 1500--1700 nm 区间，V100、A100、H100 的中位耗时约为 3.1、3.1、1.4 秒。全部使用 $R_{\rm grid}=300{,}000$，均为任何仪器卷积或重采样之前的原生网格合成耗时。APOGEE 基准只改变波长区间，不改变合成分辨率。

\begin{EN}
For the same solar atmosphere over the narrower APOGEE 1500--1700~nm interval, the V100,
A100, and H100 medians are about 3.1, 3.1, and 1.4~s.
All use $R_{\rm grid}=300{,}000$ and are native-grid synthesis timings before
any instrumental convolution or resampling. The APOGEE benchmark changes only
the wavelength interval, not the synthesis resolution.
\end{EN}

V100 运行使用 V100S 变体与 2000 条谱线的块，A100 与 H100 分别使用 4000 与 8000 条的块。V100 与 A100 耗时接近，是因为在剖面测试的太阳与 K 矮星计算中，原子与 H/He 不透明度占 80--92\%，而两者在这一主导阶段的实测时间相差不到 10\%；H100 把同一阶段约减半。这些轮廓求值与索引累积 kernel 并未使用稠密张量核矩阵乘法。各 GPU 在其可用的主机、PyTorch 与 CUDA 环境中运行，因此这些是实际应用耗时，而非孤立的架构对比。

\begin{EN}
The V100 runs use the V100S variant and 2,000-line chunks, while the A100 and
H100 runs use 4,000- and 8,000-line chunks. The V100 and A100 walls are close
because atomic and H/He opacity accounts for
80--92 percent of the profiled Sun and K-dwarf calculations, and their measured
times for this dominant stage differ by less than 10 percent. The H100 roughly
halves the same stage. These profile-evaluation and indexed-accumulation kernels
do not use dense tensor-core matrix multiplication. The GPUs ran in their
available host, PyTorch, and CUDA environments. These are therefore achieved
application timings rather than an isolated architecture comparison.
\end{EN}

合成耗时对网格采样只有弱依赖。在固定的 300--1000 nm 区间上，把 $R_{\rm grid}$ 提高 30 倍，四颗星的耗时仅增加 3.1--5.0 倍：更密的采样增加了转移列数和每条活跃轮廓内的像素数，而保留目录的扫描与每次调用的固定开销几乎不变。

\begin{EN}
Synthesis time depends only weakly on grid sampling. Over the fixed
300--1000 nm interval, increasing $R_{\rm grid}$ by a
factor of 30 raises the wall by factors of only 3.1--5.0 across the four stars.
Denser sampling adds transfer columns and pixels within each active profile,
while scanning the retained catalog and the fixed per-call floor remain nearly
unchanged.
\end{EN}

图~\ref{fig:synthesis_scaling} 检验的是波长跨度。在固定的 $R_{\rm grid}=300{,}000$ 下，窗口从 100 nm 扩到 400 nm 跨度使保留目录增大 25 倍，而 H100 耗时仅上升 1.5--1.8 倍，原 Kurucz 则上升 16--18 倍。最短窗口上 GPU 增益较小，是因为工作不足以喂饱加速器；更多波长与轮廓进入批次后，固定的单次调用开销被摊薄。此外，保留条数只是不透明度工作的上界：连续谱截断会剔除大多数“深度—谱线”对，幸存轮廓也只访问有界的像素区间。因此 H100 的加速比从 100 nm 跨度的约 3--6 倍增长到 300--1000 nm 的 36--54 倍。

\begin{EN}
Figure~\ref{fig:synthesis_scaling} instead tests wavelength span. At fixed
$R_{\rm grid}=300{,}000$, expanding the window from a 100 to a 400~nm span
increases the retained catalog by a factor of 25. The H100 wall
rises by only factors of 1.5--1.8, whereas the original Kurucz wall rises by
factors of 16--18. The smaller GPU gain for the shortest window reflects less
work to fill the accelerator. A fixed per-call floor is amortized
as more wavelengths and profiles enter the batch. Moreover, the retained count
is only an upper bound on opacity work because the continuum cutoff rejects most
depth--line pairs and surviving profiles visit bounded pixel ranges. The H100
speedup therefore grows from about $3$--$6$ for the 100~nm span to $36$--$54$
over 300--1000~nm.
\end{EN}

